% GENERATED FILE. Edit canonical lessons, then run npm run generate:books.
% canonical-source-hash: fnv1a64:58bbe3b9fb9d0655
% canonical-lessons: TA-C510-urpatti, TA-C510-poruttu, TA-C510-orunkinai, TA-C510-atukku, TA-C510-ala

\chapter{The Factory Floor}
\label{ch:ta-the-factory-floor}
\hlchaptermodality{510}

\begin{chapteropening}
\textbf{By the end of this chapter:} \emph{I can talk about making and measuring things at work.}

Complete the last lesson of chapter 510: I can talk about making and measuring things at work.
\end{chapteropening}

\section[\texorpdfstring{\ta{உற்பத்தி} (uṟpatti)}{உற்பத்தி (uṟpatti)}]{\ta{உற்பத்தி} (uṟpatti) — production; output}

\label{lesson:TA-C510-urpatti}



\begin{quote}
Before the new one: say the Tamil for to encourage, then the Tamil for an expert.
\end{quote}

\subsection*{You'll want to know: \ta{உற்பத்தி}}
\textbf{\ta{உற்பத்தி}} — \emph{uṟpatti} — \textquotedblleft{}production; output\textquotedblright{}.

\begin{quote}
\ta{இந்தத்} \ta{தொழிற்சாலையில்} \ta{உற்பத்தி} \ta{இரண்டு} \ta{மடங்கு} \ta{அதிகரித்தது}.
\end{quote}

\begin{quote}
\emph{intat toḻiṟcālaiyil uṟpatti iraṇṭu maṭaṅku atikarittatu.}
\end{quote}

\begin{quote}
Production in this factory doubled.
\end{quote}

With \ta{செய்} (\emph{cey}) it becomes a verb: \ta{உற்பத்தி} \ta{செய்}, to produce.

\begin{cousinweb}
\textbf{\ta{உற்பத்தி}} (\emph{uṟpatti}) is borrowed from Sanskrit \emph{utpatti}, \textquotedblleft{}origin, birth\textquotedblright{}, later also \textquotedblleft{}profit\textquotedblright{} (CDIAL 1812). Tamil writes the cluster \emph{tp} as \ta{ற்ப}.
\end{cousinweb}

\begin{grammarlens}[title={What you've built}]
One more word for describing what you lived and felt.
\end{grammarlens}

\subsection*{Your turn}
\begin{itemize}
\raggedright
  \item \emph{Say it:} \emph{uṟpatti}
  \item \emph{Say it:} \emph{uṟpatti}, once more
  \item \emph{Say it:} \emph{intat toḻiṟcālaiyil uṟpatti iraṇṭu maṭaṅku atikarittatu.} — Production in this factory doubled.
  \item \emph{Recall:} say the Tamil for to encourage, then the Tamil for an expert, then say \emph{uṟpatti} again
\end{itemize}

\begin{tcolorbox}[breakable,colback=teal!4,colframe=teal!35!black,title={Before you move on}]
What does \ta{உற்பத்தி} mean? (\textquotedblleft{}Production; output\textquotedblright{}.) Say it once more.
\end{tcolorbox}
\section[\texorpdfstring{\ta{பொருத்து} (poruttu)}{பொருத்து (poruttu)}]{\ta{பொருத்து} (poruttu) — to fit; to install}

\label{lesson:TA-C510-poruttu}



\begin{quote}
Before the new one: say the Tamil for an expert, then the Tamil for production.
\end{quote}

\subsection*{You'll want to know: \ta{பொருத்து}}
\textbf{\ta{பொருத்து}} — \emph{poruttu} — \textquotedblleft{}to fit; to install\textquotedblright{}.

\begin{quote}
\ta{இந்தத்} \ta{தொழிற்சாலையில்} \ta{மிதிவண்டிகளுக்குச்} \ta{சக்கரங்களைப்} \ta{பொருத்துகிறார்கள்}.
\end{quote}

\begin{quote}
\emph{intat toḻiṟcālaiyil mitivaṇṭikaḷukkuc cakkaraṅkaḷaip poruttukiṟārkaḷ.}
\end{quote}

\begin{quote}
In this factory they fit wheels to bicycles.
\end{quote}

Past (I …): \ta{பொருத்தினேன்}, \emph{poruttiṉēṉ}.

What is fitted takes -\ta{ஐ}, what it is fitted to -\ta{க்கு} or -\ta{இல்}: \ta{சுவரில்} \ta{கடிகாரத்தைப்} \ta{பொருத்து}, to put up a clock on the wall.

\begin{cousinweb}
\textbf{\ta{பொருத்து}} (\emph{poruttu}) is native Dravidian: DEDR 4541 gives \ta{பொரு} (\emph{poru}), \textquotedblleft{}to join, to unite\textquotedblright{}, and beside it \ta{பொருத்து} (\emph{poruttu}), \textquotedblleft{}to fit, to join together\textquotedblright{}.
\end{cousinweb}

\begin{grammarlens}[title={What you've built}]
One more word for describing what you lived and felt.
\end{grammarlens}

\subsection*{Your turn}
\begin{itemize}
\raggedright
  \item \emph{Say it:} \emph{poruttu}
  \item \emph{Say it:} \emph{poruttu}, once more
  \item \emph{Say it:} \emph{intat toḻiṟcālaiyil mitivaṇṭikaḷukkuc cakkaraṅkaḷaip poruttukiṟārkaḷ.} — In this factory they fit wheels to bicycles.
  \item \emph{Recall:} say the Tamil for an expert, then the Tamil for production, then say \emph{poruttu} again
\end{itemize}

\begin{tcolorbox}[breakable,colback=teal!4,colframe=teal!35!black,title={Before you move on}]
What does \ta{பொருத்து} mean? (\textquotedblleft{}To fit; to install\textquotedblright{}.) Say it once more.
\end{tcolorbox}
\section[\texorpdfstring{\ta{ஒருங்கிணை} (oruṅkiṇai)}{ஒருங்கிணை (oruṅkiṇai)}]{\ta{ஒருங்கிணை} (oruṅkiṇai) — to coordinate; to bring together}

\label{lesson:TA-C510-orunkinai}



\begin{quote}
Before the new one: say the Tamil for production, then the Tamil for to fit.
\end{quote}

\subsection*{You'll want to know: \ta{ஒருங்கிணை}}
\textbf{\ta{ஒருங்கிணை}} — \emph{oruṅkiṇai} — \textquotedblleft{}to coordinate; to bring together\textquotedblright{}.

\begin{quote}
\ta{மேலாளர்} \ta{எல்லாக்} \ta{குழுக்களையும்} \ta{ஒருங்கிணைத்தார்}.
\end{quote}

\begin{quote}
\emph{mēlāḷar ellāk kuḻukkaḷaiyum oruṅkiṇaittār.}
\end{quote}

\begin{quote}
The manager coordinated all the teams.
\end{quote}

Past (I …): \ta{ஒருங்கிணைத்தேன்}, \emph{oruṅkiṇaittēṉ}.

What is brought together takes -\ta{ஐ}: \ta{குழுக்களை} \ta{ஒருங்கிணை}.

\begin{cousinweb}
\textbf{\ta{ஒருங்கிணை}} (\emph{oruṅkiṇai}) is a native compound: \ta{ஒருங்கு} (\emph{oruṅku}), \textquotedblleft{}together\textquotedblright{}, and \ta{இணை} (\emph{iṇai}), \textquotedblleft{}to join, to unite\textquotedblright{} (DEDR 457).
\end{cousinweb}

\begin{grammarlens}[title={What you've built}]
One more word for describing what you lived and felt.
\end{grammarlens}

\subsection*{Your turn}
\begin{itemize}
\raggedright
  \item \emph{Say it:} \emph{oruṅkiṇai}
  \item \emph{Say it:} \emph{oruṅkiṇai}, once more
  \item \emph{Say it:} \emph{mēlāḷar ellāk kuḻukkaḷaiyum oruṅkiṇaittār.} — The manager coordinated all the teams.
  \item \emph{Recall:} say the Tamil for production, then the Tamil for to fit, then say \emph{oruṅkiṇai} again
\end{itemize}

\begin{tcolorbox}[breakable,colback=teal!4,colframe=teal!35!black,title={Before you move on}]
What does \ta{ஒருங்கிணை} mean? (\textquotedblleft{}To coordinate; to bring together\textquotedblright{}.) Say it once more.
\end{tcolorbox}
\section[\texorpdfstring{\ta{அடுக்கு} (aṭukku)}{அடுக்கு (aṭukku)}]{\ta{அடுக்கு} (aṭukku) — to stack; to pile up}

\label{lesson:TA-C510-atukku}



\begin{quote}
Before the new one: say the Tamil for to fit, then the Tamil for to coordinate.
\end{quote}

\subsection*{You'll want to know: \ta{அடுக்கு}}
\textbf{\ta{அடுக்கு}} — \emph{aṭukku} — \textquotedblleft{}to stack; to pile up\textquotedblright{}.

\begin{quote}
\ta{தொழிலாளர்கள்} \ta{பெட்டிகளை} \ta{ஒன்றின்} \ta{மேல்} \ta{ஒன்றாக} \ta{அடுக்கினார்கள்}.
\end{quote}

\begin{quote}
\emph{toḻilāḷarkaḷ peṭṭikaḷai oṉṟiṉ mēl oṉṟāka aṭukkiṉārkaḷ.}
\end{quote}

\begin{quote}
The workers stacked the boxes one on top of another.
\end{quote}

Past (I …): \ta{அடுக்கினேன்}, \emph{aṭukkiṉēṉ}.

The same word is a noun too: \ta{ஒரு} \ta{அடுக்கு}, a stack, a layer.

\begin{cousinweb}
\textbf{\ta{அடுக்கு}} (\emph{aṭukku}) is native Dravidian: DEDR 80 gives \ta{அடுக்கு} (\emph{aṭukku}), \textquotedblleft{}to pile one on top of another\textquotedblright{}, and the noun \textquotedblleft{}a pile, a tier\textquotedblright{}.
\end{cousinweb}

\begin{grammarlens}[title={What you've built}]
One more word for describing what you lived and felt.
\end{grammarlens}

\subsection*{Your turn}
\begin{itemize}
\raggedright
  \item \emph{Say it:} \emph{aṭukku}
  \item \emph{Say it:} \emph{aṭukku}, once more
  \item \emph{Say it:} \emph{toḻilāḷarkaḷ peṭṭikaḷai oṉṟiṉ mēl oṉṟāka aṭukkiṉārkaḷ.} — The workers stacked the boxes one on top of another.
  \item \emph{Recall:} say the Tamil for to fit, then the Tamil for to coordinate, then say \emph{aṭukku} again
\end{itemize}

\begin{tcolorbox}[breakable,colback=teal!4,colframe=teal!35!black,title={Before you move on}]
What does \ta{அடுக்கு} mean? (\textquotedblleft{}To stack; to pile up\textquotedblright{}.) Say it once more.
\end{tcolorbox}
\section[\texorpdfstring{\ta{அள} (aḷa)}{அள (aḷa)}]{\ta{அள} (aḷa) — to take measurements; to measure}

\label{lesson:TA-C510-ala}



\begin{quote}
Before the new one: say the Tamil for to coordinate, then the Tamil for to stack.
\end{quote}

\subsection*{You'll want to know: \ta{அள}}
\textbf{\ta{அள}} — \emph{aḷa} — \textquotedblleft{}to take measurements; to measure\textquotedblright{}.

\begin{quote}
\ta{தையல்காரர்} \ta{என்} \ta{கையின்} \ta{நீளத்தை} \ta{அளந்தார்}.
\end{quote}

\begin{quote}
\emph{taiyalkārar eṉ kaiyiṉ nīḷattai aḷantār.}
\end{quote}

\begin{quote}
The tailor measured the length of my arm.
\end{quote}

Past (I …): \ta{அளந்தேன்}, \emph{aḷantēṉ}.

The stem ends in -\ta{அ}, so the past is in -\ta{ந்த}-: \ta{அளந்தேன்}. The noun \ta{அளவு} (\emph{aḷavu}), a size, comes from it.

\begin{cousinweb}
\textbf{\ta{அள}} (\emph{aḷa}) is native Dravidian: DEDR 295 gives \ta{அள} (\emph{aḷa}), \textquotedblleft{}to measure, to limit\textquotedblright{}, and its noun \ta{அளவு} (\emph{aḷavu}), \textquotedblleft{}a measure, a size\textquotedblright{}.
\end{cousinweb}

\begin{grammarlens}[title={What you've built}]
One more word for describing what you lived and felt.
\end{grammarlens}

\subsection*{Your turn}
\begin{itemize}
\raggedright
  \item \emph{Say it:} \emph{aḷa}
  \item \emph{Say it:} \emph{aḷa}, once more
  \item \emph{Say it:} \emph{taiyalkārar eṉ kaiyiṉ nīḷattai aḷantār.} — The tailor measured the length of my arm.
  \item \emph{Recall:} say the Tamil for to coordinate, then the Tamil for to stack, then say \emph{aḷa} again
\end{itemize}

\begin{tcolorbox}[breakable,colback=teal!4,colframe=teal!35!black,title={Before you move on}]
What does \ta{அள} mean? (\textquotedblleft{}To take measurements; to measure\textquotedblright{}.) Say it once more.
\end{tcolorbox}
